\section{Discusión}

La red de asociaciones proteicas del fenotipo
HP:0012347 presentó una densidad de 0.1293
y una transitividad global de 0.6263.

Estos resultados muestran una conectividad global
limitada junto con una tendencia al agrupamiento
local. Esta organización es compatible con la
existencia de comunidades, identificadas
posteriormente mediante el algoritmo de Louvain.
Los clusters recuperan complejos y rutas conocidos: 
las enzimas \textit{ALG} de la síntesis del oligosacárido ligado
a dolicol, el complejo COG, que participa en el tráfico dentro
del Golgi, y las subunidades de la V-ATPasa. Que los genes más
centrales por centralidad armónica sean \textit{ALG1}, \textit{ALG3}, \textit{ALG6}, 
\textit{PMM2} y \textit{ALG11} es coherente con el papel de esta ruta como paso previo a 
la N-glicosilación.

El enriquecimiento global confirma que el fenotipo está asociado a la glicosilación, 
pero este resultado era esperable, porque los genes se eligieron precisamente por su
relación con ella. Más informativo es el análisis por clusters, que separa la síntesis
del precursor (cluster 1), el metabolismo del dolicol y la manosilación (cluster 2) 
y el tráfico intracelular de proteínas (cluster 4). En cambio, el cluster 3,
formado por subunidades de la V-ATPasa, solo devolvió el término genérico 
\textit{biological regulation}, y el cluster 2 solo términos muy amplios de metabolismo 
del fosfato. Estos resultados no deben sobreinterpretarse, pues la identidad de esos 
clusters se deduce de los genes que los componen y no del enriquecimiento.

El análisis tiene varias limitaciones. Nueve de los 50 genes 
(\textit{B4GALT1}, \textit{SLC35A2}, \textit{UGGT1}, \textit{GORAB}, \textit{MAN2B2},
\textit{MGAT2}, \textit{SLC37A4}, \textit{SLC39A8} y \textit{MAN1B1}) no tienen 
interacciones con puntuación mayor o igual que 800 y quedan fuera de la red y 
de todas las métricas. Bajar el umbral los incorporaría a costa de interacciones 
menos fiables. Además, el algoritmo de Louvain es estocástico, por lo que el número y 
la composición de los clusters pueden variar entre ejecuciones y versiones de 
\textit{igraph}. Por último, el fondo de solo 41 genes reduce la potencia del ORA por 
cluster, y solo 22 de los 50 genes tienen anotado el término 
\textit{protein N-linked glycosylation} en GO, lo que apunta a que las anotaciones de GO 
y de la HPO son incompletas.

\subsection{Organización topológica y heterogeneidad de la red}

El análisis topológico de la red permitió ampliar la
caracterización de los genes asociados al fenotipo
HP:0012347 más allá de la identificación de comunidades
y del enriquecimiento funcional.

La red presentó una densidad del 12.93\% y una
transitividad global de 0.6263. La coexistencia de
una densidad limitada y una transitividad relativamente
elevada sugiere una organización en grupos locales
de proteínas interconectadas, en consonancia con las
comunidades identificadas mediante el algoritmo de Louvain.

El análisis de la distribución del grado mostró una
desviación típica de 2.7559 y un coeficiente de variación
del 53.30\%, lo que refleja una heterogeneidad apreciable
en la conectividad de las proteínas.

Los genes \textit{ALG1} y \textit{ALG3} destacaron
por presentar 11 asociaciones cada uno, superando
el doble del grado medio de la red. Ambos pertenecen
al cluster 1, identificado previamente mediante Louvain,
que presentó un enriquecimiento significativo en
procesos relacionados con la N-glicosilación de proteínas.

Asimismo, otros genes con elevada conectividad,
como \textit{ALG6}, \textit{PMM2}, \textit{DPM1}
y \textit{DPAGT1}, pertenecen a los dos primeros
clusters de la red.

Estos resultados sugieren que las proteínas con mayor
grado se concentran en comunidades relacionadas con
los mecanismos de glicosilación y su metabolismo
asociado.

No obstante, una elevada conectividad en STRING no
implica necesariamente una mayor esencialidad biológica,
relevancia clínica o relación causal con la gravedad
del fenotipo. Por ello, resulta conveniente complementar
el análisis del grado con otras medidas de centralidad
y robustez de la red.


\subsection{Relevancia estructural y conectividad
entre comunidades}

El análisis de centralidades permitió distinguir
diferentes formas de importancia estructural dentro
de la red de asociaciones proteicas.

Aunque ALG1 y ALG3 presentaron el mayor grado,
ALG2 destacó por su centralidad de intermediación
y por constituir un punto de articulación.

Este resultado sugiere que ALG2 desempeña un papel
relevante en la conectividad de la red que no puede
explicarse exclusivamente por el número de
asociaciones directas.

Asimismo, ATP6V0A2 presentó una elevada
intermediación y constituyó otro punto de
articulación, a pesar de no encontrarse entre
los genes con mayor grado.

El análisis de participación permitió identificar
a MPI, SRD5A3 y PMM2 como candidatos relevantes
para estudiar las relaciones entre comunidades.

PMM2 destacó simultáneamente por su conectividad,
intermediación y participación, lo que sugiere
un papel estructural relevante tanto a nivel local
como en la organización global de la red.

La identificación de seis puntos de articulación
pone de manifiesto la existencia de nodos cuya
eliminación puede modificar la conectividad de
la red.

Por otro lado, la correlación moderada entre grado
e intermediación justifica la utilización de
múltiples medidas de centralidad, puesto que
cada una caracteriza diferentes propiedades
estructurales.

No obstante, los resultados deben interpretarse
teniendo en cuenta que STRING representa asociaciones
funcionales y que la selección del umbral de
confianza condiciona la topología obtenida.

La relevancia topológica de una proteína no implica
necesariamente su esencialidad biológica ni permite
establecer directamente relaciones causales con
la gravedad del fenotipo.

Estos hallazgos constituyen una base para evaluar
posteriormente la vulnerabilidad y robustez de la red
mediante simulaciones de eliminación de nodos.
